\begin{proposition}
\label{proposition-galois}
\begin{reference}
This is a weak version of \cite[Expos\'e V]{SGA1}.
The proof is borrowed from \cite[Theorem 7.2.5]{BS}.
\end{reference}
Let $(\mathcal{C}, F)$ be a Galois category. Let $G = \text{Aut}(F)$
be as in Example \ref{example-from-C-F-to-G-sets}. The functor
$F : \mathcal{C} \to \textit{Finite-}G\textit{-Sets}$
(\ref{equation-remember}) an equivalence.
\end{proposition}

\begin{proof}
We will use the results of Lemma \ref{lemma-epi-mono} without further mention.
In particular we know the functor is faithful.
By Lemma \ref{lemma-tame} we know that for any connected $X$ the
action of $G$ on $F(X)$ is transitive. Hence for general $X$ the functor
$F$ sends the connected components of $X$ (which exist by
an axiom of a Galois category) to the $G$-orbits of $F(X)$.
Let $X$ and $Y$ be objects and let
$s : F(X) \to F(Y)$ be a map. Then the graph
$\Gamma_s \subset F(X) \times F(Y)$ of $s$
is a union of connected components. Hence there exists a
union of connected components $Z$ of $X \times Y$,
which comes equipped with a monomorphism $Z \to X \times Y$,
with $F(Z) = \Gamma_s$. Since $F(Z) \to F(X)$ is bijective
we see that $Z \to X$ is an isomorphism and we conclude
that $s = F(f)$ where $f : X \cong Z \to Y$ is the composition.
Hence $F$ is fully faithful.

\medskip\noindent
To finish the proof we show that $F$ is essentially surjective.
It suffices to show that $G/H$ is in the essential image for
any open subgroup $H \subset G$ of finite index.
By definition of the topology on $G$ there exists a finite
collection of objects $X_i$ such that
$$
\Ker(G \longrightarrow \prod\nolimits_i \text{Aut}(F(X_i)))
$$
is contained in $H$. We may assume $X_i$ is connected
for all $i$. We can choose a Galois object $Y$ mapping
to a connected component of $\prod X_i$ using
Lemma \ref{lemma-galois}. Choose an isomorphism $F(Y) = G/U$
in $G\textit{-sets}$ for some open subgroup $U \subset G$.
As $Y$ is Galois, the group
$\text{Aut}(Y) = \text{Aut}_{G\textit{-Sets}}(G/U)$ acts transitively
on $F(Y) = G/U$. This implies that $U$ is normal. Since
$F(Y)$ surjects onto $F(X_i)$ for each $i$ we see that
$U \subset H$. Let $M \subset \text{Aut}(Y)$ be the finite subgroup
corresponding to
$$
(H/U)^{opp} \subset (G/U)^{opp} = \text{Aut}_{G\textit{-Sets}}(G/U)
= \text{Aut}(Y).
$$
Set $X = Y/M$, i.e., $X$ is the coequalizer
of the arrows $m : Y \to Y$, $m \in M$.
Since $F$ is exact we see that $F(X) = G/H$ and the
proof is complete.
\end{proof}